\documentclass[11pt]{article}
\usepackage[margin=0.9in]{geometry}
\usepackage{fontspec}
\usepackage{amsmath,amssymb,amsthm}
\usepackage{xurl}
\begin{document}\sloppy
\begin{center}
{\Large GROK H0203: cube 19 of \texttt{a40ff54e580f5d0a} at $L=9$ is 132 of 210 subcubes UNSAT with no gap, and 185 of cube 0's subcubes are certified}\\[4pt]
2026-10-06 02:03 America/Los\_Angeles \quad branch \texttt{main}, tip \texttt{cde425a1} \quad automation 41774ea2\\[2pt]
\textbf{P versus NP: OPEN.}
\end{center}

\section*{Peer delta since H0145 (\texttt{cde425a1})}
Checked 02:04\,PT. No commits landed after \texttt{cde425a1}, so nothing is new under \texttt{research/main/} or \texttt{research/agents/grok/}. Root \texttt{main.tex} is not touched by this note.

\section*{1. Literature}
None needed. The next claim is about our own finite search.

\section*{2. Claims recorded}
\textbf{Claim DRAT185 (COMPUTATIONAL).} 185 of cube 0's 210 subcubes carry drat-trim-verified refutations. The 4 new ones are 134, 129, 202 and 154. That leaves 25; the batch is now on subcube 87.

\textbf{Partial status of cube 19 (COMPUTATIONAL, solver-trusted, incomplete).} First-step cube 19, $(0,0,4,14)$, has 132 of its 210 subcubes UNSAT under CaDiCaL and none SAT. It is \emph{not} closed. Every index up to the largest finished one (133) is either finished or running (127 and 129 are running).

\textbf{Status of first-step cubes (COMPUTATIONAL, solver-trusted).} Unchanged from H0145: 48 of 75 kept cubes closed, 27 open, 0 SAT.

\section*{3. Smallest honest test}
\begin{itemize}
\item \textbf{Gap control} (inline Python over \texttt{a40\_L9\_split.out} plus the process list, 02:04\,PT). Cube 19 has finished indices with maximum 133; the only missing indices at or below 133 are 127 and 129, and both have live solver processes, as do 134, 135 and 136.
\item \textbf{Map control} (\texttt{/workspace/scratch\_h0033\_1006/experiments/cube\_check.py}, rerun 02:03\,PT). The cube-19 lines seen so far have distinct indices and tuples, each equal to cube 0's step-2 tuple for that index, all UNSAT.
\item \textbf{Cap control} (inline Python, 02:03\,PT). Across all 1391 UNSAT split lines there is no \texttt{TIMEOUT\_or\_exit} line, no SAT line, and nothing over 500\,s; the slowest is still 376.5\,s against the 600\,s cap.
\item \textbf{Recount} (\texttt{/workspace/scratch\_h2333\_1005/experiments/recount.py}, rerun 02:03\,PT). 81 earlier verified subcubes plus 104 batch lines, all \texttt{s VERIFIED}, in queue order, with 104 distinct CNF hashes, give 185.
\end{itemize}

\section*{4. What this changes}
The certified count moves from 181 to 185. The closed count stays 48 of 75. The bounds stay $9\le L(\texttt{a40ff54e580f5d0a},2)\le 11$ (COMPUTATIONAL); $L=9$ itself stays open while 27 first-step cubes are open. The largest confirmed price stays one.

\section*{Effect on the standing candidate}
No claim is falsified and no charged resource improves. The time-state-representation surface stays CANDIDATE. The SAT-to-charged-model compiler is OPEN, and so is the product-span obstruction. These results concern one fixed function and imply no lower bound for any family.

\section*{Next smallest step}
Let phase B finish cube 19 and rerun the gap and map controls before calling it closed. Let the DRAT batch finish its last 25.

\bigskip
\noindent\textbf{Status sentence.} P versus NP remains OPEN: the SAT-to-charged-model compiler and the product-span obstruction are open load-bearing arrows. This pass records that 185 of cube 0's 210 subcubes of \texttt{a40ff54e580f5d0a} at tight $L=9$ carry drat-trim-verified refutations (COMPUTATIONAL) and that first-step cube 19 is 132 of 210 subcubes UNSAT with no gap but not yet closed; 48 of 75 kept first-step cubes are closed, 27 open, 0 SAT, and $9\le L(f,2)\le 11$ is unchanged. No claim is falsified, no charged resource improves, and this is not a proof of anything about P versus NP.
\end{document}
